% 黄道十二宫（综述）
% license CCBYSA3
% type Wiki

本文根据 CC-BY-SA 协议转载翻译自维基百科\href{https://en.wikipedia.org/wiki/Astrological_sign}{相关文章}。

\begin{figure}[ht]
\centering
\includegraphics[width=6cm]{./figures/e66f90583cecd88b.png}
\caption{出自《贝里公爵的豪华时祷书》的解剖学黄道人} \label{fig_Astsi_1}
\end{figure}
在西方占星术中，占星星座即黄道，被划分为十二个各占30度的区段，这些区段由太阳在地球天空中所观测到的360度轨道路径所穿越。星座的顺序从春季的第一天开始，即被称为“白羊宫起点”的位置，也就是春分点。占星星座依次为白羊座、金牛座、双子座、巨蟹座、狮子座、处女座、天秤座、天蝎座、人马座、摩羯座、宝瓶座和双鱼座。西方黄道起源于巴比伦占星术，后来受到希腊化文化的影响。每一个星座均以太阳每年穿越天空时经过的一个星座命名。这一观测在简化且流行的太阳星座占星术中被尤为强调。几个世纪以来，由于地球自转轴进动$^{[1]}$，西方占星术中的黄道分区逐渐与其所命名的星座发生偏离，而印度占星术的测量方法则对这种偏移进行了修正。$^{[2]}$ 占星术（即一种基于天象的预兆体系）在中国和西藏文化中也有所发展，但这些占星体系并不以黄道为基础，而是涉及整个天空。

占星术是一种伪科学。$^{[3]}$ 对其理论基础$^{[4]}$以及对相关主张的实验验证$^{[5]}$的科学研究表明，它既不具有科学有效性，也缺乏解释力。对于人格特征与出生月份之间表面相关性的现象，存在更为合理的解释，例如人类季节性出生所产生的影响。

根据占星学的观点，天体现象与人类活动之间遵循“上如其下”的原则，因此星座被认为代表着具有特征性的表达模式。$^{[6]}$ 直到19世纪之前，科学天文学与西方占星术一样，均使用相同的黄道分区。

目前，不同的占星体系采用了多种测量和划分天空的方法，尽管黄道星座的名称与符号传统在很大程度上保持一致。西方占星术以春分点和至点（即热带年中昼夜等长、最长与最短的时刻）为基准进行测量，而印度占星术则沿着赤道平面进行测量，对应的是恒星年。
\subsection{西方黄道星座}  
\subsubsection{历史}
